% !TeX root = ../papers/nc-p3.tex
\section{第二長方形の反対階数と階数十四の特異K3曲面}
\label{sec:second-rectangle-opposite-ranks}
\begin{theorem}\label{thm:second-rectangle-opposite-ranks}
台143から147の水平表示で $b+c\ne0$ とする。$A=u+1/u$, $B=v+1/v$, $P=z+w$, $Q=zw$ と置き
\[
 f=(A+B)(P-2)-(B+2)(Q-1),\quad
 g=(A+B)(P+2)+(A+2)(Q-1)
\]
とする。反対階数は $u=v=1$ で8、ちょうど一方が1なら10、残りでは
$16-\mathbf1_{f=0}-\mathbf1_{g=0}$ である。
二つの階数15因子の閉包は有理三次元多様体で、一位置移動で交換される。
その交叉は階数14の滑らかな既約曲面である。指定した座標でのその極小滑らか射影模型は
Picard数20、Néron--Severi判別式の絶対値16の楕円K3曲面である。
台147でのこの曲面の一般の無印同型類は有理曲面でパラメータ付けられる。
階数14の曲面と両方の階数15因子には、全点が標準Koszul左右AS正則な非空開集合がある。
\end{theorem}

\begin{proof}[階数と分岐]
$C=(u+v)(1+uv)$, $H=u(v+1)^2$, $G=v(u+1)^2$ とする。
四色の反対行列の行列式から定理\ref{thm:horizontal-rectangle-low-fibers}の $M$ の単元を除くと、
順に $(v-1)^2f,(u-1)^2g,(u-1)^2(v-1)^2,(u-1)^2(v-1)^2$ である。
全色に単元の二次小行列式がある。三次小行列式のイデアルはそれぞれ
$(v-1),(u-1),((u-1)(v-1)),((u-1)(v-1))$ となる。
最初の色で共通因子を除いた後には、単元倍を除き $z-1,w-1$ が得られ、第二色では $z+1,w+1$ が得られる。
両方の同時消滅は $Q\ne1$ で排除される。残る二色では共通因子の単元倍が一つの小行列式で得られる。
$u=1$ なら $uv f=-(v+1)^2(z-1)(w-1)$ は可逆、$v=1$ なら
$uv g=(u+1)^2(z+1)(w+1)$ は可逆である。これで全点での階数が従う。
単独の二因子ではそれぞれ $w$ を一次で解け、係数と定数項が互いに素である。
一位置移動は $(A,B,P,Q)\mapsto(B,A,-P,Q)$ と作用する。

$N=A+B+4$ と置くと $f=g=0$ で $N(Q-1)=-4(A+B)$ なので $N\ne0$ であり、
\[
 P=2(A-B)/N,\qquad Q=1-4(A+B)/N .
\]
従って $(z,w)$ はこれらを和と積に持つ順序対である。
$Y=N(z-w)/4$ とすれば
\[
 Y^2=D(A,B)=A^2+AB+B^2+2A+2B-4.
\]
基底のトーラスで分岐曲線の特異点があるとすれば $u=\pm1$ または $v=\pm1$ である。
実際 $D=D_A=D_B=0$ に解はなく、$D(2,B)=(B+2)^2$, $D(-2,B)=B^2-4$ である。
このうち $(u,v)=(1,-1),(-1,1)$ は元のトーラス開集合の外、$(-1,-1)$ は $N=0$ である。
よって階数14の層は滑らかで、二因子はそこで横断的に交わる。
どちらかの座標が1の面上では前記の可逆性により両因子に同時には特殊化しない。
\end{proof}

\begin{proof}[楕円曲面と無印商]
変数 $A$ を底とし $v^2D(A,v+1/v)$ を四次式として取ると
\[
 (vY)^2=v^4+(A+2)v^3+(A^2+2A-2)v^2+(A+2)v+1 .
\]
$v=0$ に有理切断を持つ。二元四次式の不変式 $I,J$ はそれぞれ $A$ の次数4,6、
判別式は $(A-2)^3(A+2)^3(3A+10)^2$ である。
$A=2$ で $I$ は単純零点、$-2,-10/3$ と無限遠の極小表示では単元である。
従って四つのファイバーは $2,-2,-10/3,\infty$ の順に
$\mathrm{III},I_3,I_2,I_4$、基本直線束が次数1の有理楕円面である。
$A=u+1/u$ は $2,-2$ で分岐し、基底変換後のファイバーは
\[
\begin{array}{c|rrrr}
 u&1&-1&-3,-1/3&0,\infty\\\hline
 \text{型}&I_0^*&I_6&I_2,I_2&I_4,I_4
\end{array}
\]
となる。判別式次数24、切断と標準束公式から楕円K3となる。
有理楕円面の自明格子の根の階数は7なのでMordell--Weil階数は1である。
四根を逆数対に分ける分割が有理2捩れ点を与え、加法的ファイバーの成分群により捩れはちょうど位数2である。
有理楕円面の unimodular なNéron--Severi格子から原始生成元の高さは $4/(2\cdot3\cdot2\cdot4)=1/12$ である。

K3の自明格子の階数は19で、引き戻した非捩れ切断によりPicard数は20となる。
四次式の判別式は $u$ の関数体で平方である。四根の三分割に対する作用は、逆数対の分割を固定し
偶置換だけから来るので全ての分割を固定する。従って全2捩れが有理であり、$I_0^*$ の成分群から捩れの位数は4である。
引き戻した切断の高さは $1/6$。ファイバー寄与の分母は全て12を割るため、
これが非捩れ切断の $m$ 倍と捩れの和であれば $1/(6m^2)\in\frac1{12}\mathbb Z$ である。
よって $|m|=1$。格子の判別式の絶対値は
$4\cdot6\cdot2^2\cdot4^2\cdot(1/6)/4^2=16$ となる。

一般点での文字枠の復元は下記の剰余検査から従う。従って指定始点の無印商には
$I:(u,v,Y)\mapsto(1/u,1/v,-Y)$ だけを追加して割ればよい。
コンパクト化上の $u=v=-1,Y=0$ では $D_A=D_B=-4$ であり、局所方程式の二次部が非退化な通常二重点である。
対合の接作用は三座標で全て $-1$ なので解消の例外円錐曲線を点ごとに固定する。
この解消はK3の極小模型にも残る（例えば二重被覆のコンパクト化は分岐の型 $(4,4)$ を持つ。
四隅は二本の異なる接線を持つ節点で、$(1,-1),(-1,1)$ の残る特異性は
分岐曲線の二重接触、二重被覆の $A_3$ であり、全解消はcrepantである）。
固定曲線を持つ対合は正則二形式に $-1$ で作用する。その商は滑らかで、不規則数は0、
二重被覆公式で反標準因子の二倍は非零な有効分岐因子と線型同値である。
従って小平次元が負となり、商は有理曲面である。
\end{proof}

\begin{proposition}\label{prop:second-rectangle-rank-fourteen-period}
上の階数14曲面で $z=w$ の層は滑らかな有理曲線で、一般点は標準正則かつ最小周期2である。
その外の曲面の一般点は最小周期4である。
\end{proposition}
\begin{proof}
$l=(u-1)/(u+1),m=(v-1)/(v+1),r=lm$ と置く。$D=0$ は単元を除き
$1-3l^2m^2+l^2m^2(l^2+m^2)=0$ である。
\[
 S=r(l+m)/(r+1),\qquad E=r(l-m)/(r-1)
\]
と置くと $S^2=2r-1$, $E^2=-2r-1$ となる。逆に滑らかな円錐曲線 $S^2+E^2=-2$ から
$r=(S^2+1)/2$, $l=((r+1)S+(r-1)E)/(2r)$, $m=((r+1)S-(r-1)E)/(2r)$,
$z=w=-SE$ を回収できる。除いた分母の零集合は曲線の成分を含まない。
$J^2$ は指定始点の同型 $I$ を除けば $z,w$ を交換し、一般に奇数移動は固定点を持たない。

階数14曲面とその対角曲線で、それぞれ $u=2$ を取る。前者は $v=3$,
$59z^2+10z-81=0$, $w=-10/59-z$、後者は
$4v^4+18v^3+37v^2+18v+4=0$, $z=w=(A-B)/N$ である。
素数1009の単純根 $z=883$, $v=729$ で台147の全52係数は非零、各位置の隣接二位置に対する
二次モーメント一つだけで異なる四固有値を持つ。
全始点で五次までの経路次元と四次中心の次元 $(3,2,0,0)$ を確認し、中心四式の商は
長さ13まで $(1+T+T^2+T^3)^4$ の係数を持つ。中心核の特殊化については、
主成分一般点からの次元下界と剰余点の上界が一致し、経路も五次まで標準であるため単元小行列で持ち上がる。
前節までの中心商の上界と一般点の全次数下界との比較で、この二数体の点で正則性を得る。
四つの一般中心式の有限判定および枠の四固有値はその各点の近傍で開く。
これにより曲面、含む両因子、対角曲線の各非空開集合での主張が従う。
\path{computations/verify_second_rectangle_opposite_ranks.py} が行列式・曲面・円錐曲線を、
\path{computations/verify_second_rectangle_rank_fourteen_samples.py} が二つの数体の検査を再現する。
\end{proof}
